% =============================================================================
% Chapter 2 --- Background
% =============================================================================

\chapter{Background}
\label{ch:background}

This chapter establishes only the context required to justify the
controller's closed-loop architecture: the energetic case for in-memory
computing (IMC), the physical mechanism by which a crossbar implements a
matrix--vector multiplication (MVM), the Y-Flash device that this controller
drives, and the write--verify problem the device's non-idealities impose on
the surrounding digital logic.

\section{In-Memory Computing and the MVM Kernel}
\label{sec:bg-imc}

Modern data-centric workloads --- ANN inference foremost among them --- are
bottlenecked by the energy cost of moving operands between memory and
arithmetic units rather than by the arithmetic itself; this is the
\emph{memory wall}~\cite{Sebastian2020IMC, Ielmini2018RRAM, Verma2019IMC}.
IMC sidesteps the bottleneck by embedding the computation \emph{inside} the
memory array. For the dominant ANN kernel, MVM over a stored weight
matrix, the crossbar topology is structurally aligned with the dataflow:
weights are stored at the crosspoints, input activations are applied on
the rows, and the per-cell contributions sum physically on the columns.

In algebraic form the array implements $\vec{y} = W\vec{x}$ through two
physical laws. \textbf{Ohm's law} at each crosspoint produces a current
$I_{ij} = G_{ij}V_i = w_{ij}x_i$, where the programmable conductance
$G_{ij}$ encodes the synaptic weight $w_{ij}$. \textbf{Kirchhoff's current
law} on each column sums these currents into
$I_j = \sum_i w_{ij}x_i$, yielding one MAC per column per
read~\cite{Sebastian2020IMC, Verma2019IMC, Ielmini2018RRAM}.
\Cref{fig:mvm-crossbar-array} sketches the canonical arrangement.

The energy advantage of this organization is intrinsic: an entire MVM row
is consumed in a single access rather than as one memory cycle per weight.
The benefit, however, is delivered only when the surrounding periphery ---
addressing, pulse sequencing, calibration, error
handling~\cite{Sebastian2020IMC} --- preserves the array-level efficiency.
Three requirements follow directly for the digital controller: every cell
must be individually \emph{addressable}, every operation must be
\emph{pulse-shaped} on the appropriate terminals, and every program must
be \emph{verifiable} against a target value. The architecture in
\Cref{ch:architecture} maps onto these three requirements.

\begin{figure}[h]
    \centering
    \includegraphics[width=0.85\textwidth]{example-image.png}
    \caption{$M \times N$ memristor crossbar performing an analog
        matrix--vector multiplication. Row read-voltages $\{V_i\}$ drive
        programmable conductances $G_{ij}$; per-column currents
        $I_j = \sum_i G_{ij} V_i$ are digitized at the array periphery.}
    \label{fig:mvm-crossbar-array}
\end{figure}

\section{The Y-Flash Floating-Gate Memristor}
\label{sec:bg-yflash}

Among the candidate IMC devices --- RRAM, PCM, MRAM, and floating-gate
variants~\cite{Ielmini2018RRAM, Wang2022YFlash} --- the analog core this
controller drives uses the \emph{Y-Flash} floating-gate memristor, the
device technology under development in the supervisor's
group~\cite{Wang2022YFlash}.

The cell is built in a standard single-poly CMOS process without
additional masks and contains two transistors --- a longer-channel
\emph{read} transistor and a shorter-channel \emph{injection} transistor
--- that share a common floating gate and a common drain. Three operations
are exposed through distinct external bias configurations: \textbf{Read}
at $V_R < 2.5\,\text{V}$ exposes the stored charge as a subthreshold
current in the \SI{1}{\nano\ampere}--\SI{5}{\micro\ampere} range;
\textbf{Program} at $V_P > 4\,\text{V}$ injects hot electrons onto the
floating gate, raising threshold and reducing conductance; \textbf{Erase}
at $V_E > 7\,\text{V}$ discharges the gate via band-to-band
tunneling~\cite{Wang2022YFlash}. The selection is performed at the cell
itself --- de-biasing the source defeats programming on un-selected
cells --- which makes the device naturally compatible with per-cell
one-hot row/column selection from the digital side. The physical
arrangement of the two-transistor cell and the shared floating gate is
sketched in \Cref{fig:yflash-structure}.

\begin{figure}[h]
    \centering
    \includegraphics[width=0.7\textwidth]{yflash_structure.png}
    \caption{Physical structure of the Y-Flash floating-gate memristor.
        The cell pairs a longer-channel \emph{read} transistor with a
        shorter-channel \emph{injection} transistor that share a single
        floating gate and a common drain, and is fabricated in a
        standard single-poly CMOS process without additional masks. The
        three operations --- Read at low gate bias, Program by
        channel-hot-electron injection onto the floating gate, and Erase
        by band-to-band tunneling --- are selected entirely by the
        external bias configuration, while per-cell selection is
        achieved at the source terminal so that un-selected cells in the
        same row or column remain undisturbed.}
    \label{fig:yflash-structure}
\end{figure}

Two properties of the device dictate the structure of the digital
controller. First, all three operations rely on tunneling or hot-carrier
injection, whose effective amplitude depends \emph{exponentially} on
applied voltage and pulse duration: the controller is therefore
responsible for accurate pulse-time shaping, not just bias selection.
Second, the analog conductance is tunable across a continuum of
intermediate states with low cycle-to-cycle variability relative to
filamentary alternatives~\cite{Wang2022YFlash}, but the post-programming
conductance is still displaced from its target by an amount that depends
on local fabrication state, charge history, and pulse shape. These
displacements are what motivate the closed-loop control policy.

\section{The Write--Verify Problem}
\label{sec:bg-writeverify}

No single unconditional program pulse can be relied upon to place a
resistive cell at its intended conductance. Sebastian et~al.\ describe
the canonical mitigation as an iterative \emph{write--read--decide}
strategy: each program pulse is interleaved with a verify read, the
readback is compared against the target, and subsequent pulses are
adapted (or replaced by an erase) until the cell converges or a retry
budget is exhausted~\cite{Sebastian2020IMC, Ielmini2018RRAM,
Wang2022YFlash}. Three architectural consequences for the controller
follow:

\begin{itemize}
    \item \textbf{One-hot addressing at the boundary.} The host
          transaction selects exactly one cell per block through
          concatenated one-hot fields, mirroring the per-cell deselect
          mechanism the Y-Flash exposes during program and
          erase~\cite{Wang2022YFlash}.
    \item \textbf{A multi-bit pulse-mode bus.} Read, program, erase, and
          inference each map to a distinct time-multiplexed burst encoding
          rather than a generic write bit-pattern, matching the
          device-level distinction between the three bias
          configurations.
    \item \textbf{A hierarchical FSM with explicit recovery branches.}
          The controller must implement a program--verify--(re-program or
          erase) loop with a bounded retry budget, so that a single
          mis-programmed cell cannot bring down the system: it is either
          re-programmed, erased and re-attempted, or declared a
          failure~\cite{Ielmini2018RRAM, Sebastian2020IMC,
          Wang2022YFlash}.
\end{itemize}

These three consequences are the design contract of the digital
controller. The remainder of the report describes the RTL that implements
that contract.
